% ============================================================================
% CAPÍTULO 6 - SEGURANÇA, GUARDRAILS E EXPLICABILIDADE
% ============================================================================

\chapter{Segurança, Guardrails e Explicabilidade}\label{cap:seguranca}

A aplicação de \acrshort{ia} generativa em contexto clínico exige camadas de proteção que vão além da acurácia do modelo. Este capítulo detalha os três pilares de confiabilidade implementados no assistente: \textbf{guardrails de segurança}, \textbf{explicabilidade} e \textbf{auditoria}.

\section{Princípios de Segurança em IA Médica}

O design de segurança do assistente segue três princípios fundamentais:

\begin{enumerate}
    \item \textbf{Nunca substituir o médico:} O sistema é uma ferramenta de \emph{apoio à decisão}, não um tomador de decisão. Toda resposta deve ser validada por um profissional de saúde.

    \item \textbf{Falhar de forma segura:} Quando houver incerteza, o sistema prefere recusar ou alertar a arriscar uma informação potencialmente danosa. Uma resposta conservadora com baixa confiança é preferível a uma resposta confiante mas imprecisa.

    \item \textbf{Transparência total:} Cada resposta deve ser rastreável --- quais fontes foram consultadas, qual o nível de confiança, quais regras foram verificadas e se houve retries.
\end{enumerate}

Estes princípios se materializam nos três subsistemas descritos a seguir.

\section{Guardrails: Validação Automática de Respostas}\label{sec:guardrails}

O agente de guardrails (\texttt{src/flows/agents/guardrails.py}) aplica quatro regras de validação a cada resposta gerada pelo agente de raciocínio. A validação é \textbf{determinística} --- baseada em regex e palavras-chave --- o que garante previsibilidade e auditabilidade, sem depender de um segundo modelo de linguagem.

\subsection{Regra 1: Proibição de Prescrição Direta}

O assistente não pode prescrever medicamentos. Frases como ``prescrevo Losartana 50mg'' ou ``tome 2 comprimidos de 8/8h'' são detectadas e rejeitadas.

\begin{lstlisting}[caption={Padrões regex para detecção de prescrição direta.}]
PRESCRIPTION_PATTERNS = [
    r"prescrev[oa]",
    r"tome\s+\d+",
    r"administr[ea]\s+\d+\s*mg",
    r"iniciar?\s+(com\s+)?\d+\s*mg",
    r"receita[r:]",
    r"usar?\s+\d+\s*(mg|ml|comprimido)",
]
\end{lstlisting}

Quando detectado, o feedback instrui o modelo a reformular como \emph{sugestão} e indicar validação médica. Por exemplo, ``considerar o uso de Losartana 50mg, conforme protocolo, sujeito à avaliação do médico responsável''.

\subsection{Regra 2: Validação Humana Obrigatória}

Toda resposta deve conter um \textbf{disclaimer} indicando necessidade de validação por profissional de saúde. A verificação busca palavras-chave como ``validação'', ``profissional de saúde'', ``médico responsável'' ou ``avaliação clínica''.

\begin{lstlisting}[caption={Palavras-chave de disclaimer de validação humana.}]
DISCLAIMER_KEYWORDS = [
    "validação", "validar", "profissional de saúde",
    "médico responsável", "avaliação clínica",
    "decisão clínica", "confirmar com",
]
\end{lstlisting}

Se ausente, o guardrails reprova a resposta com feedback solicitando a inclusão da ressalva.

\subsection{Regra 3: Proibição de Diagnóstico Definitivo}

O assistente não pode afirmar diagnósticos de forma categórica. Frases como ``o paciente tem diabetes'' ou ``confirmo que é pneumonia'' são detectadas e rejeitadas.

\begin{lstlisting}[caption={Padrões regex para detecção de diagnóstico definitivo.}]
DEFINITIVE_DIAGNOSIS_PATTERNS = [
    r"o paciente tem\b",
    r"o diagnóstico é\b",
    r"trata-se de\b",
    r"confirmo que\b",
    r"certamente (é|tem)\b",
]
\end{lstlisting}

O feedback orienta o uso de linguagem probabilística: ``sugere'', ``é compatível com'', ``deve ser investigado''.

\subsection{Regra 4: Consistência com Contexto}

Uma verificação heurística avalia se a resposta é compatível com o volume de contexto disponível. Se protocolos e/ou dados do paciente foram recuperados mas a resposta tem menos de 50 caracteres, o guardrails considera que o modelo pode ter ``alucinado'' uma resposta genérica e solicita elaboração.

\subsection{Execução Sequencial das Regras}

As quatro regras são executadas sequencialmente. Se qualquer regra falha, a resposta é marcada como \texttt{reprovado} e todos os feedbacks são concatenados para fornecer ao modelo de raciocínio uma instrução completa de correção.

\begin{lstlisting}[caption={Pipeline de validação dos guardrails.}]
GUARDRAIL_CHECKS = [
    ("no_prescription", _check_no_prescription),
    ("human_validation", _check_human_validation),
    ("no_definitive_diagnosis", _check_no_definitive_diagnosis),
]

for check_name, check_fn in GUARDRAIL_CHECKS:
    passed, message = check_fn(draft_response)
    if not passed:
        all_passed = False
        feedback_parts.append(message)

# Consistencia com contexto (separada)
passed, message = _check_context_consistency(
    draft_response, protocols, patient_data
)
\end{lstlisting}

\subsection{Mecanismo de Retry com Feedback}

Conforme descrito na Seção~\ref{sec:agentes} do Capítulo~\ref{cap:multiagente}, respostas reprovadas são devolvidas ao agente de raciocínio com o feedback concatenado. O mecanismo permite até \textbf{2 retries}. Se o limite for atingido, a resposta segue para a explicabilidade com advertências adicionais --- o sistema nunca bloqueia completamente o fluxo.

\begin{importantbox}[Decisão de Design: Por que não bloquear?]
Um bloqueio total geraria frustração no uso clínico e poderia atrasar decisões urgentes. A abordagem de ``seguir com warnings'' permite que o médico receba a informação acompanhada de alertas claros, preservando a autonomia do profissional.
\end{importantbox}

O feedback injetado no prompt de retry tem formato específico:

\begin{lstlisting}[caption={Feedback de retry injetado no prompt do raciocínio.}]
ATENCAO - SUA RESPOSTA ANTERIOR FOI REPROVADA PELO GUARDRAILS:
Motivo: Prescricao direta detectada: 'prescrevo'.
Reformule como sugestao e indique validacao medica.
| Resposta nao inclui disclaimer de validacao humana.
Adicione uma nota indicando que a resposta deve ser
validada por um profissional de saude.
Corrija os problemas apontados nesta nova resposta.
\end{lstlisting}

\section{Explicabilidade}\label{sec:explicabilidade}

O agente de explicabilidade (\texttt{src/flows/agents/explicabilidade.py}) transforma o rascunho validado em uma resposta \textbf{transparente}, adicionando quatro componentes: citações, confiança, advertências e disclaimer.

\subsection{Citações e Rastreabilidade}

Toda resposta inclui a lista de fontes consultadas, extraídas automaticamente dos metadados dos protocolos recuperados pelo ChromaDB e dos dados do paciente. As fontes são \textbf{separadas por relevância}:

\begin{itemize}
    \item \textbf{Fontes diretas} --- diretamente relacionadas ao quadro do paciente:
    \begin{itemize}
        \item Protocolos cujo conteúdo contém termos das comorbidades: \texttt{[Pathway-Diabetes, Seção: Tratamento]}
        \item Dados do paciente: \texttt{[Prontuário: Maria Santos]}, \texttt{[Exame: HbA1C, Data: 2026-01-15]}
    \end{itemize}
    \item \textbf{Fontes complementares} --- informação geral que pode ser útil, mas sem relação direta com as comorbidades do paciente.
\end{itemize}

A deduplicação é feita via conjunto (\texttt{set}), garantindo que cada citação apareça apenas uma vez mesmo quando múltiplos chunks vêm do mesmo documento.

\subsection{Avaliação de Confiança}

A confiança é calculada por um sistema de pontuação que considera a \textbf{quantidade} e \textbf{relevância} do contexto recuperado, a \textbf{disponibilidade} de dados do paciente e eventuais \textbf{penalizações} por retries:

\begin{table}[H]
\centering
\caption{Sistema de pontuação para avaliação de confiança.}
\label{tab:confianca_scoring}
\rowcolors{2}{fiap-light}{white}
\begin{tabular}{l r l}
\toprule
\textbf{Fator} & \textbf{Pontos} & \textbf{Justificativa} \\
\midrule
3+ protocolos recuperados & $+3$ & Múltiplas fontes convergentes \\
1--2 protocolos recuperados & $+2$ & Ao menos uma fonte de evidência \\
Distância média $< 0.5$ & $+2$ & Alta relevância semântica \\
Distância média $< 1.0$ & $+1$ & Relevância moderada \\
Dados do paciente disponíveis & $+2$ & Contexto individualizado \\
Cada retry do guardrails & $-1$ & Indica problemas na geração \\
\bottomrule
\end{tabular}
\end{table}

A classificação final segue três faixas:

\begin{itemize}
    \item \textbf{Alta} ($\geq 5$ pontos): múltiplas fontes relevantes e contexto completo;
    \item \textbf{Média} ($\geq 3$ pontos): alguma fundamentação disponível;
    \item \textbf{Baixa} ($< 3$ pontos): informação insuficiente para alta confiança.
\end{itemize}

\subsection{Advertências Contextuais}

O sistema gera advertências automáticas quando detecta situações que merecem atenção:

\begin{itemize}
    \item \textbf{Confiança baixa:} ``Poucos dados disponíveis para fundamentar a resposta.'';
    \item \textbf{Sem protocolos:} ``Nenhum protocolo clínico foi encontrado para esta consulta.'';
    \item \textbf{Retries realizados:} ``Resposta foi refinada N vezes pelo guardrails de segurança.''
\end{itemize}

\subsection{Disclaimer Obrigatório}

Independentemente do conteúdo, toda resposta inclui o disclaimer:

\begin{warningbox}[Disclaimer do Assistente]
Esta resposta é uma sugestão baseada em protocolos clínicos e dados disponíveis. Toda decisão clínica deve ser validada pelo médico responsável.
\end{warningbox}

\subsection{Formato da Resposta Final}

A resposta final apresentada ao médico segue uma estrutura padronizada:

\begin{lstlisting}[caption={Formato da resposta final com explicabilidade.},language={}]
[Resposta clinica do agente de raciocinio]

Fontes consultadas:
  . [Pathway-Diabetes, Secao: Tratamento]
  . [Pathway-Hipertensao, Secao: Exames de Seguimento]
  . [Exame: HbA1C, Data: 2026-01-15]

Fontes complementares:
  . [Pathway-PacienteCirurgico, Secao: Exames Complementares]

Confianca: Alta

Advertencias:
  . Resposta foi refinada 1x pelo guardrails de seguranca.

Esta resposta e uma sugestao baseada em protocolos clinicos
e dados disponiveis. Toda decisao clinica deve ser validada
pelo medico responsavel.
\end{lstlisting}

\section{Auditoria e Logging}\label{sec:auditoria}

O agente de logger (\texttt{src/flows/agents/logger\_agent.py}) é o nó final em todos os fluxos do grafo. Sua responsabilidade é \textbf{consolidar} o estado completo da execução em um registro estruturado e \textbf{persistir} para consulta posterior.

\subsection{Estrutura do Registro de Auditoria}

Cada interação gera um registro JSON com seis seções:

\begin{table}[H]
\centering
\caption{Seções do registro de auditoria.}
\label{tab:audit_sections}
\rowcolors{2}{fiap-light}{white}
\begin{tabularx}{\textwidth}{l X}
\toprule
\textbf{Seção} & \textbf{Conteúdo} \\
\midrule
\texttt{session} & Query original e ID do paciente \\
\texttt{triagem} & Tipo de query classificado e entidades extraídas \\
\texttt{context} & Quantidade de protocolos, fontes consultadas, disponibilidade de dados \\
\texttt{response} & Tamanho do rascunho vs.\ final, confiança, fontes citadas, advertências \\
\texttt{guardrails} & Resultado (aprovado/reprovado), feedback, número de retries \\
\texttt{agent\_trace} & Lista cronológica de todos os agentes executados com timestamps \\
\bottomrule
\end{tabularx}
\end{table}

A seção \texttt{agent\_trace} é construída de forma \textbf{acumulativa}: cada agente adiciona sua entrada ao campo \texttt{audit\_log} do estado. Ao final, o logger consolida todas as entradas em um registro único, criando um \textbf{trace completo} da execução.

\subsection{Persistência em JSONL}

Os registros são persistidos em arquivos diários no formato \textbf{JSONL} (JSON Lines):

\begin{lstlisting}[style=bash,caption={Estrutura de persistência dos logs de auditoria.}]
logs/
  audit_2026-03-06.jsonl   # Um registro JSON por linha
  audit_2026-03-07.jsonl
  ...
\end{lstlisting}

O formato JSONL foi escolhido por três razões:

\begin{enumerate}
    \item \textbf{Append-only:} Novos registros são adicionados ao final do arquivo sem reescrever os anteriores, minimizando risco de corrupção.
    \item \textbf{Streaming:} Cada linha é um JSON válido independente, facilitando processamento incremental com ferramentas como \texttt{jq} ou ingestão em sistemas de monitoramento.
    \item \textbf{Simplicidade:} Não requer infraestrutura adicional (banco de dados, message broker), adequado para o escopo do projeto.
\end{enumerate}

\subsection{Rastreabilidade End-to-End}

Com o registro de auditoria, é possível reconstruir o caminho completo de qualquer consulta:

\begin{enumerate}
    \item Qual foi a query original e como foi classificada;
    \item Quais protocolos foram recuperados e de quais fontes;
    \item Se e quais dados do paciente foram consultados;
    \item Qual foi o rascunho gerado pela LLM;
    \item Se o guardrails aprovou, reprovou ou realizou retries;
    \item Qual o nível de confiança e quais advertências foram emitidas;
    \item O timestamp de cada etapa para análise de latência.
\end{enumerate}

Esta rastreabilidade é essencial para:

\begin{itemize}
    \item \textbf{Depuração:} Identificar por que uma resposta foi inadequada;
    \item \textbf{Compliance:} Demonstrar que o sistema seguiu todas as regras de segurança;
    \item \textbf{Melhoria contínua:} Analisar padrões de reprovação para refinar regras ou prompts.
\end{itemize}

\section{Tratamento de Queries Fora de Escopo}

Quando a triagem classifica uma consulta como \texttt{fora\_de\_escopo}, o fluxo é encurtado: a query vai direto ao logger, que gera uma resposta padrão educada:

\begin{infobox}[Resposta para Query Fora de Escopo]
``Desculpe, esta pergunta está fora do escopo do assistente médico. Posso ajudar com consultas sobre protocolos clínicos, dados de pacientes, tratamentos e condutas médicas.''
\end{infobox}

O registro de auditoria é gerado normalmente, permitindo rastrear quais tipos de perguntas estão fora do escopo --- informação útil para expandir a cobertura do sistema no futuro.

\section{Limitações e Trabalhos Futuros}

As implementações atuais de segurança, embora funcionais, possuem limitações conhecidas:

\begin{enumerate}
    \item \textbf{Guardrails baseados em regex:} Padrões fixos podem gerar falsos positivos (ex: ``o paciente tem histórico de\ldots'' é flagrado pela regra de diagnóstico definitivo) ou falsos negativos (prescrições reformuladas de forma indireta). Uma evolução natural seria utilizar um classificador baseado em LLM como segunda camada de validação.

    \item \textbf{Confiança heurística:} O sistema de pontuação é determinístico e não considera a \emph{qualidade} do conteúdo dos protocolos, apenas quantidade e proximidade semântica. Métricas de faithfulness (fidelidade ao contexto) poderiam enriquecer a avaliação.

    \item \textbf{Logging local:} A persistência em JSONL local é adequada para demonstração, mas em produção seria necessário integrar com sistemas centralizados de observabilidade (ex: ELK Stack, Datadog, OpenTelemetry).

    \item \textbf{Consentimento e privacidade:} O sistema registra a query completa no log de auditoria. Em ambiente de produção com dados reais, seria necessário implementar anonimização automática e políticas de retenção de dados conforme LGPD.
\end{enumerate}
